% =====================================================================
%  GMN3 §9  $A_1$ 理论与叶层（$A_1$ theories and laminations）
%  原文：Lamination-Paper.tex 5051--5359 行；PDF p.79--85（印刷 p.78--84）
%  编号公式 (9.1)--(9.11) 共 11 条，全部在本块内，故 \eqref 可用；
%  图 6 幅（9.1--9.6），图名照抄原文 \insfig 第一参数：
%    dipole / path-passage / path-deform / path-deform-2 /
%    limit-laminations / limit-laminations-higher
%  跨块引用（式 (6.14)、§10、附录 F.1/F.2）暂写原文编号并加 % TODO，
%  合稿时统一改 \ref/\eqref
% =====================================================================

\section{$A_1$ 理论与叶层}
\label{sec:A1}

本节与第 10 节% TODO 合稿改 \ref{sec:Examples}
专门讨论 $A_1$ 理论。在这类理论中，可以给出分解 (6.14)% TODO 合稿改 \eqref{eq:y1}
的显式例子——该分解把线算子的 vev 用 $\tCY_\gamma$ 表示出来——从而算出带框 BPS 简并度。本节先给出一些必要的预备知识；具体结果见第 10 节。% TODO 合稿改 \ref{sec:Examples}

\subsection{技术回顾}

我们首先回顾文献 \cite{Gaiotto:2009hg} 中用到的若干技术，更多细节请参阅该文。

固定一点 $u \in \CB$。$u$ 对应 $C$ 上一个亚纯二次微分 $\phi_2$，而 $\phi_2$ 决定了 SW（Seiberg--Witten）曲线 $\Sigma$，其形式为 $\lambda^2 = \phi_2$。电荷格为 $\Gamma = H_1(\Sigma, \IZ)^{odd}$。$\phi_2$ 在 $C$ 上一个正则奇点（比如 $z = 0$）附近的行为由一个边界条件决定：
\begin{equation}
 \phi_2 \sim m^2 \frac{dz^{\otimes 2}}{z^2}.
\end{equation}
这里 $m$ 是对应于理论的一个相关形变的质量参数。在非正则奇点处，$\phi_2$ 有高阶极点。对一般的 $u$，$\phi_2$ 只有单零点；$\CB^{sing} \subset \CB$ 是 $\phi_2$ 的零点不全为单零点的轨迹。

在 $C$ 上任何既非 $\phi_2$ 的极点、也非其零点的点附近，我们可以用 $w = \int \sqrt{\phi_2}$ 定义一个局部坐标（差一个符号）。我们把 $C$ 上这样一条曲线称为「相位为 $\vartheta$ 的 WKB 曲线」：它映到 $w$ 平面后成为一条倾角为 $\vartheta$ 的直线。以 $t$ 参数化的一般 WKB 曲线，当 $t \to \pm\infty$ 时渐近于 $\phi_2$ 的极点。在满足 $m e^{i\vartheta} \notin i\IR$ 的正则奇点附近，WKB 曲线描写盘旋落入该奇点的指数螺线。在 $\phi_2$ 有 $k+2$ 阶极点\footnote{在文献 \cite{Gaiotto:2009hg} 中，非正则奇点的阶是用一个整数 $L$ 表示的；但那个记号会与我们表示线算子的记号冲突。}
的非正则奇点附近，WKB 曲线聚成 $k$ 个不同的族，每一族渐近于 $k$ 条「WKB 射线」之一。

在文献 \cite{Gaiotto:2009hg} 中，我们还引入了 $C$ 的一个三角剖分 $T_\WKB(u, \vartheta)$——更确切地说，是如下曲面的三角剖分：从 $C$ 出发，在每个非正则奇点周围挖去一个小圆盘，再在边界上标出一些点，把每个边界圆分成 $k$ 段弧。当不致引起混淆时，我们滥用记号，把这个带标记的曲面也称作 $C$。$T_\WKB(u, \vartheta)$ 的边是一般的 WKB 曲线。顶点是 WKB 曲线的渐近端点：每个正则奇点是一个顶点；每个非正则奇点给出一组 $k$ 个顶点，每段弧上各一个，此外还有这些顶点与斯托克斯射线之间的一个对应。$T_\WKB(u, \vartheta)$ 还被一些额外的离散数据「装饰」，如下所述。

正如我们已经提到的，在其复结构 $J^\zeta$ 下，$\CM$ 是带奇点的平坦联络的模空间。给定这样一个联络，再给定 $C$ 上一条 WKB 曲线，我们在每个渐近无穷远处都得到一条特别的「小平坦截面」（在相差一个重标度的意义下定义）：它就是当我们沿 WKB 曲线走进奇点时呈指数衰减的那条截面。在正则奇点处，这条小平坦截面不依赖于我们考虑的是哪一条 WKB 曲线；而且它是绕该奇点的单值变换的两个本征截面之一，相应的本征值为
\begin{equation}
 \mu = \exp \left[ \pm 2 \pi i \left( R \zeta^{-1} m - 2 m^{(3)} - R \zeta \bar{m} \right) \right].
\end{equation}
在 $\phi_2$ 有 $k+2$ 阶极点的非正则奇点处，情况稍复杂一些，原因在于斯托克斯现象：如我们所说，WKB 曲线聚成 $k$ 个族，每一族渐近于 $k$ 条 WKB 射线之一，而每条 WKB 射线携带一条不同的小平坦截面。所以无论哪种情形，$T_\WKB$ 的每个顶点都附有一条特别的平坦截面。

利用装饰三角剖分 $T_\WKB(u, \vartheta = \arg \zeta)$，我们在文献 \cite{Gaiotto:2009hg} 中定义了一族 $\CM$ 上的典范坐标函数 $\CX_\gamma$；其定义见附录 F.1。% TODO 合稿改 \ref{app:fgcoords}
我们还在文献 \cite{Gaiotto:2009hg} 中定义了 $\Gamma$ 上一个典范的二次加细 $\sigma$，于是可以再定义一族坐标函数 $\CY_\gamma = \sigma(\gamma) \CX_\gamma$。文献 \cite{Gaiotto:2009hg} 的论证表明：对 $\gamma \in \Gamma$，这些 $\CY_\gamma$ 确实与本文一直使用的那些 $\CY_\gamma$ 一致。另一方面，上面讨论带框 BPS 简并度时，我们还需要属于某个扩张格 $\Gamma_\CL$ 的电荷 $\gamma$ 所对应的 $\CY_\gamma$。对这些电荷，文献 \cite{Gaiotto:2009hg} 没有给出 $\CY_\gamma$ 的定义。由于 $2\gamma \in \Gamma$，$\CY_\gamma^2$ 是有定义的，所以困难在于确定一个平方根的符号。这里我们不给出确定这些符号的处方，而是改在扭空间 $\tCM_\CL$ 上工作，以绕开这个困难：在这个空间上符号困难不那么严重，我们确实可以对所有 $\gamma \in \Gamma_\CL$ 定义所需的函数 $\tCY_\gamma$。这件事在附录 F.2 中完成。% TODO 合稿改 \ref{app:traffic}

最后，在文献 \cite{Gaiotto:2009hg} 的附录 A 中，我们回顾了一个把 $\CM$ 上的和乐函数 $\Tr \Hol_\wp \CA$ 用坐标 $\CY_\gamma$ 展开的配方，但带有若干我们一般没有固定的符号歧义。在附录 F.2% TODO 合稿改 \ref{app:traffic}
中我们对这个配方稍作改进：给出 $\tCM$ 上的扭和乐函数用 $\tCY_\gamma$ 的展开，且不带任何符号歧义。我们将用这个展开来提取带框 BPS 简并度 $\fro(u, L_\zeta, \gamma; y = 1)$。

\subsection{缺失线算子之谜}
\label{sec:missing}

正如我们刚刚描述的，$\CS$ 类中的一个 $A_1$ 理论由如下选择刻画：一条黎曼面 $C$，其上带有若干「正则」或「非正则」穿刺点。迄今关于这类理论的大部分文献都集中在正则穿刺点上，但一些最自然的（也是组合上最简单的）例子其实需要非正则穿刺点：例如纯 $SU(2)$ 理论这样的渐近自由理论总是需要非正则穿刺点，文献 \cite{Gaiotto:2009hg} 所述 Argyres--Douglas CFT 的实现也是如此。

在这样的理论中如何构造线算子？我们可以试着沿用上面的做法、并照文献 \cite{Drukker:2009tz} 所述，由 $C$ 上的闭路径得到它们。然而稍加考察例子就会看出，这样并不能给出我们预期存在的全部线算子。例如，纯 $SU(2)$ 理论对应于带两个非正则穿刺点的 $C = \IC\IP^1$ \cite{Gaiotto:2009hg}；于是按同痕分类，$C$ 上只有一条非自交闭路径。这一条可以等同于 Wilson 线算子。可是 Wilson--'t Hooft 算子在哪里？

稍后我们就会看到，缺失的那些线算子对应于某些\emph{整叶层}（integral lamination），即允许以非正则穿刺点为端点的路径之组合。

\subsection{去耦味}
\label{subsec:Decouple-Flavor}

我们来考察当共形理论退化为渐近自由理论时，线算子会发生什么变化。

我们从一个用曲线 $C'$ 定义的 $A_1$ 理论 $T'$ 出发。设在 $C'$ 上某个局部坐标为 $z$ 的片内有两个正则奇点。我们设想这两个奇点靠得很近，比如位于 $z = \pm \eps$，质量参数为 $m_1, m_2$。于是有
\begin{equation}
\begin{split}
\phi_2 & = \left( \frac{m_1^2}{(z-\eps)^2} + \frac{m_2^2}{(z+\eps)^2} + \cdots \right) \, dz^{\otimes 2}\\
& = \frac{F(z)}{(z - \eps)^2 (z + \eps)^2} \, dz^{\otimes 2}.
\end{split}
\end{equation}

考虑把 $F(z)$ 在 $z = 0$ 附近展开成泰勒级数。通过调节 $+\eps$、$-\eps$ 两个奇点处的留数 $m_1$、$m_2$，我们可以任意固定这个级数的常数项与线性项。现在设 $\eps \to 0$ 而保持这两项有限。在这个极限下可以写
\begin{equation} \label{eq:limit-q}
\phi_2 = \left(\frac{\Lambda}{z^2} + \frac{m}{z} + \cdots\right)^2  \, dz^{\otimes 2}
= \left( \frac{\Lambda^2}{z^4} + \frac{2 \Lambda m}{z^3} + \cdots\right)  \, dz^{\otimes 2}
\end{equation}
其中 $\eps^2(m_1^2 + m_2^2) \to \Lambda^2$，且 $\frac{(m_1^2 - m_2^2)^2}{m_1^2 + m_2^2} \to m^2$。特别地，$m_1$ 与 $m_2$ 都按 $1/\eps$ 标度，而 $m_1 + m_2$ 与 $m_1 - m_2$ 中总有一个是 $1$ 阶的。我们还可以进一步取 $m \to \infty$、$\Lambda \to 0$ 而保持 $\Lambda m$ 固定的极限。

于是在这个标度极限下，原理论 $T'(m_1, m_2)$「退化」为一个理论 $T(m, \Lambda)$。新理论用一条带非正则奇点的曲线 $C$ 定义：事实上，\eqref{eq:limit-q} 正是「有两条斯托克斯射线从中伸出」的非正则奇点所对应的形式。

现在我们可以讨论极限理论 $T$ 中的线算子了：从理论 $T'$ 的线算子出发，观察它们在极限下变成什么。对一些线算子，答案相当平凡：任何不经过 $C'$ 上那两个相撞奇点附近的闭路径 $\wp'$，都简单地对应于 $C$ 上的一条闭路径 $\wp$，而和乐 $\Hol_{\wp'} \CA(\zeta)$ 极限到 $\Hol_\wp \CA(\zeta)$。因此，很自然地可以延拓文献 \cite{Drukker:2009tz}：猜想 $\wp$ 对应于理论 $T'$ 中的一个线算子 $L_\wp$，并写下
\begin{equation}
 L_{\wp'} \to L_\wp
\end{equation}
在极限下成立。

更有意思的情形是 $\wp'$ 确实从两个相撞奇点之间穿过，以致在 $\eps \to 0$ 的极限下被「夹断」。这时我们得多费些力气。首先，局部分析表明：在相撞奇点附近，当 $\eps$ 很小、$\vartheta$ 取一般值时，WKB 叶状结构看起来就像磁偶极子的场，如图 \ref{fig:dipole} 所示。
\insfig{dipole}{$C'$ 上一对相邻正则奇点的小邻域内的 WKB 叶状结构。}
特别地，我们看到从相撞奇点邻域中长出两束特别的叶。这两束叶渐近地趋于两条特别的射线。而且，取 $\eps \to 0$ 后对非正则奇点附近作 WKB 近似可以看出，这两条特别的射线恰好就是从该奇点伸出的 WKB 射线。（比较文献 \cite{Gaiotto:2009hg} 的图 42，那里画的非正则奇点有 $4$ 条 WKB 射线伸出。）类似地，沿这些 WKB 曲线的小平坦截面的空间在这个极限下光滑地相互关联。由此可以这样选择小平坦截面 $s_1$、$s_2$ 作为 $\eps$ 的函数的归一化：使它们极限到两条斯托克斯射线各自的那条小平坦截面。

现在，研究叶层 $\wp'$ 的一个方便技巧是「插入一组完备态」，把 $\wp'$ 切成若干段。给定一段路径 $q$ 以及沿 $q$ 的两条平坦截面 $s_1$、$s_2$，沿 $q$ 的一般平坦截面 $s$ 可以重写成两项之和：
\begin{equation}
 s = \frac{s_1 \wedge s}{s_1 \wedge s_2} s_2 + \frac{s \wedge s_2}{s_1 \wedge s_2} s_1.
\end{equation}
于是，沿 $q$ 从起点 $i$ 到终点 $f$ 的平行移动算子同样是两项之和：
\begin{equation}  \label{eq:two-terms}
 s(f) = \frac{s_1(i) \wedge s(i)}{s_1 \wedge s_2} s_2(f) + \frac{s(i) \wedge s_2(i)}{s_1 \wedge s_2} s_1(f).
\end{equation}

我们将把它用于如下情形：取 $q$ 为 $\wp'$ 的一小段，并取 $s_1$、$s_2$ 为两个奇点附近的小平坦截面，沿最显然的最短路线延拓到 $q$。见图 \ref{fig:path-passage}。
\insfig{path-passage}{$\wp'$ 的子路径 $q$。}

我们的目标是把答案改写成不涉及任何在奇点之间穿行的路径的形式。为此，我们分别处理 \eqref{eq:two-terms} 中的两项。先看第一项。我们按图 \ref{fig:path-deform} 所示对 $\wp'$ 作形变。
\insfig{path-deform}{对 $\wp'$ 的一种形变，使 $q$ 的两个端点能在 $\eps \to 0$ 极限下存活。}
经过这个形变，式中出现的量 $s_1(i)$ 与 $s_2(f)$ 在两个奇点相撞的极限下仍然有意义。另一方面，当我们试图计算 $s_1 \wedge s_2$ 时却遇到一个困难：$s_1(i)$ 与 $s_2(f)$ 有限这一事实在这里帮不上忙，因为我们必须在\emph{同一个}点上计算 $s_1$ 与 $s_2$。这就要把其中之一沿 $q$ 搬运过去，而 $q$ 正是在两个奇点之间穿行。为了绕开这个困难，我们可以把 $q$ 重新表示成一个绕 $z_1$ 的小回路与一条不会被夹断的新路径 $q'$ 的复合，如图 \ref{fig:path-deform-2} 所示。
\insfig{path-deform-2}{对 $\wp'$ 的进一步形变。}
沿 $q$ 计算 $s_1 \wedge s_2$，于是给出一个因子 $\mu_1^{-1}$（来自 $s_1$ 绕 $z_1$ 的单值变换——回忆我们把 $\mu_i$ 定义为 $s_i$ 绕 $z_i$ 的逆时针单值变换）乘上一个有限的部分。这个有限的部分仍是 $s_1 \wedge s_2$，只不过现在应理解为沿 $q'$ 而非 $q$ 搬运到公共点后的结果。对 \eqref{eq:two-terms} 的第二项可以作十分类似的分析。于是平行移动变成
\begin{equation}
 s(f) = \mu_1^{-1} \frac{s_1(i) \wedge s(i)}{s_1 \wedge s_2} s_2(f) + \mu_2 \frac{s(i) \wedge s_2(i)}{s_1 \wedge s_2} s_1(f),
\end{equation}
其中当 $\eps \to 0$ 时，除 $\mu_1$ 与 $\mu_2$ 外所有因子都有限且非零。

现在考虑 $\wp'$ 上平行移动在 $SL(2,\IC)$ 基础表示中的\emph{迹}。我们已经把沿 $q$ 的搬运写成了两个秩 1 算子之和。相应地，迹分解为两个不同对象的形式和：
\begin{equation} \label{eq:lp-decomp}
 L_{\wp'} = \mu_1^{-1} L_1 + \mu_2 L_2.
\end{equation}
我们把 $L_{1,2}$ 中的每一个表示成若干有向路径\emph{段}之并：这些段的端点落在绕非正则奇点 $\CP$ 的圆 $S^1(\CP)$ 上，并携带权重 $\pm 1$。见图 \ref{fig:limit-laminations}。
\insfig{limit-laminations}{极限对象 $L_{1,2}$。}
每条有向路径段的两端都落在从非正则奇点伸出的斯托克斯扇区上，于是在两端 $i$、$f$ 处各有特别的小平坦截面 $s_i$、$s_f$。$L_{1,2}$ 中任何一个的 vev 都是逐段之积：一条从 $i$ 到 $f$、携带权重 $k$ 的段贡献 $(s_i \wedge s_f)^k$。

在严格的 $\eps \to 0$ 极限下，对 $m_{1,2}$ 与 $\zeta$ 的相位的一般取值，$\mu_1$ 与 $\mu_2$ 都趋于 $0$ 或 $\infty$。\eqref{eq:lp-decomp} 中的两项会有一项压过另一项。因此可以说，在 $\eps \to 0$ 极限下 $L_{p'}$ 趋近于 $L_1$、$L_2$ 两个对象之一，再乘上一个无穷大的「重正化」因子。我们把 $L_1$ 与 $L_2$ 视为理论 $T'$ 的良好线算子。

于是 $\eps \to 0$ 极限的效果，是以一种特定的方式把路径 $\wp'$「切断」。更一般地，让多对奇点相撞，就可以把 $\wp'$ 切成许多不同的段。

上面考虑的是基础表示中的和乐迹。我们同样可以考虑更高的自旋表示，比如基础表示的 $k$ 次对称幂。那种情形下 \eqref{eq:lp-decomp} 的类比是一个分解
\begin{equation}
 L_{\wp', k} = \sum_{n=0}^k \mu_1^{-k-n} \mu_2^{n} L_1^{k-n} L_2^n.
\end{equation}
\insfig{limit-laminations-higher}{极限对象 $L_1^{k-n} L_2^n$。}
对象 $L_1^{k-n} L_2^n$ 画在图 \ref{fig:limit-laminations-higher} 中；正如记号所暗示的，它的 vev 为 $\inprod{L_1^{k-n} L_2^n} = \inprod{L_1}^{k-n} \inprod{L_2}^{n}$。在 $\eps \to 0$ 极限下，这个和同样由 $n = 0$ 或 $n = k$ 那一项主导。

我们上面的构造与文献 \cite{fgtoappear} 中给出的一个构造密切相关。

\subsection{叶层}
\label{subsec:Lams-Def}

上面我们发现：在带有「从中伸出两条斯托克斯线」的非正则奇点的 $A_1$ 理论中，良好线算子的类不仅包括 $C$ 上携带 $SL(2,\IC)$ 表示的闭路径，还包括终止于非正则奇点周围斯托克斯扇区、携带整数权重的开路径。我们的推导用到了某种取极限的程序，但很自然地会相信：即便非正则奇点不是由这种程序产生的，附着在开路径上的线算子仍然有意义——毕竟，在 $(2,0)$ 理论中定义表面算子的物理在 $C$ 上是局域的。

如果 $C$ 带有更一般的非正则奇点、斯托克斯线的条数任意，那会怎样？有一个自然的方案，它推广了我们上面针对两条斯托克斯线所得的结果，并且可以由多个正则奇点相撞推导出来。\footnote{这个推导的完整细节我们尚未完成。}
它与 Fock--Goncharov 在一个密切相关的语境中用过的构造非常相似 \cite{MR2349682}。沿用（但略有滥用）他们的术语，我们把所考虑的对象称为\emph{叶层}（lamination，是「整 $\CA_0$-叶层」的简称）。

回忆一下：当 $C$ 有非正则穿刺点时，我们挖去每个非正则穿刺点周围的小圆盘，并把每个边界圆分成若干段弧。$C$ 上的一个叶层 $L$ 是 $C'$ 上一族曲线 $\wp_i$ 之并：每条曲线或者闭合，或者在边界弧上有端点；它们不自交、也互不相交，并按同痕分类。每条闭曲线携带 $SL(2,\IC)$ 的一个非平凡不可约表示，但为下文的技术方便有一个例外：绕单个正则穿刺点的闭曲线携带一个整数权重（可正可负），而不携带表示。此外，每条开曲线 $\wp_i$ 携带一个整数权重 $k_i$，要求为正，但有一个例外：若一条开曲线可以收缩到边界上恰含一个标记点的弧段，则其权重允许为负。终止在边界每一段弧上的所有曲线的权重之和必须为零。还有如下等价关系：若一条曲线可缩，或可收缩到边界上不含标记点的弧段，则它可以被去掉；若叶层 $L$ 含有两条同痕的曲线 $\wp_i$、$\wp_j$，权重分别为 $k_i$、$k_j$，则它等价于一个新的叶层：去掉 $\wp_j$，而 $\wp_i$ 携带权重 $k_i + k_j$。

\emph{我们猜想：对任何 $A_1$ 理论，叶层与单线算子一一对应。}

附着在这样一个叶层上的线算子，其真空期望值对叶层的各分量是乘性的。附着在闭路径 $\wp$ 上的因子照例是沿 $\wp$ 的和乐之迹——但绕单个穿刺点、权重为 $k$ 的闭路径除外，对它们我们附着因子 $\mu^k$。对权重为 $k$ 的开路径，因子为
\begin{equation} \label{eq:open-vev}
 (s_{1} \wedge s_{2})^{k},
\end{equation}
其中 $s_1$、$s_2$ 分别是该路径起点与终点顶点处的小平坦截面。虽然附着在开路径上的因子依赖于小平坦截面 $s_i$ 的归一化选择，但当我们在叶层中对各分量取乘积时，这一归一化会抵消。这正是我们在上面的定义中对权重 $k_i$ 加上限制的原因。

上面我们没有仔细固定 vev 的整体符号。事实上，虽然有确定这个符号的处方，但下文中我们并不需要它。当我们确实要留意符号时，我们改在扭局部系统的空间上工作；在后一个空间上有一种自然的定符号方式，见附录 F.2。% TODO 合稿改 \ref{app:traffic}
